\section{Introduction}
\label{sec:intro}

% Four moves, roughly one paragraph each. Do not exceed ~900 words for the
% capstone version; journals will want ~700.

% =====================================================================
% REVIEW 2026-09-01 . Work through these, then delete the block.
%
% [R1] BLOCKING, FACTUAL. Move 1 says "land change between 2020 and 2024".
%      Your snapshots are 2012, 2016, 2020, 2024 (Decision Log D03). A
%      reader who catches this in paragraph one reads everything after it
%      with suspicion. Fix first.
% [R2] "but also today" : you have no 2026 snapshot. Latest observed year is
%      2024 and forecasts start at 2030. Either cut "today" or say what data
%      would support it.
% [R3] SCAFFOLD IS STALE, not your prose. Move 4 below says "State the three
%      objectives" and points at "the Hypotheses sheet". There is no
%      Hypotheses tab in the workbook any more, and the Decision Log now
%      carries a fourth strand: D26 predictive extension to 2030/2040/2050,
%      D27 hindcast validation, D28 conservation prioritisation with
%      Protect/Defend/Restore/Mitigate labels. Update Move 4 to match.
% [R4] PROTOCOL DRIFT, worth resolving before you write Move 4. The protocol
%      was git-tagged locked on 18 Aug. D26-D28 expand scope after that. The
%      forecasting work is legitimate and predeclared in the Decision Log,
%      but the locked protocol and the manuscript currently describe
%      different studies. Add an amendment-log entry in protocol section 19
%      so the sequence is on the record.
% [R5] No citations in the drafted prose. The scaffold asks for the range
%      assessment and the global decline synthesis. Nothing in the
%      Literature Log covers the bottleneck dating or the range-decline
%      claim yet, so those need sourcing before they can stay.
% =====================================================================

\subsection*{Move 1: Why cheetah connectivity, and why outside protected areas}

\TODO{Global decline and range contraction. The fact that most southern African
cheetahs persist on unprotected land is the whole reason a connectivity paper is a
conservation paper rather than a park-management paper. Cite the range assessment
and the global decline synthesis. Keep to context; no resistance calibration here.}

The southern African cheetah has had their range substantially reduced from their original range, with their current extant range being on majority unprotected land. Cheetah range on protected land is governed with park-management initiatives and are much easier to identify conservation efforts. Unprotected land serves a different type of challenge, with many stakeholders and many variables that are not easily controllable. This paper not only looks at land change between 2020 and 2024, but also today, and provides a predictive look at what future land permeability for cheetah movement looks like based on the trends identified.

% --- notes on the paragraph above ---
% [a] "reduced from their original range" : repeats "range", and "original"
%     is undated. Original when?
% [b] "current extant range" : prohibited phrase, and redundant. See
%     abstract R3.
% [c] BROKEN SENTENCE. "Cheetah range on protected land is governed with
%     park-management initiatives and are much easier to identify
%     conservation efforts." Subject and verb disagree, "governed with" ->
%     "governed by", and the final clause has no grammatical object. I can
%     guess you mean protected land has a single manager and a clearer
%     intervention path, but a committee should not have to guess. Rewrite
%     from scratch rather than patching.
% [d] "serves a different type of challenge" -> "poses".
% [e] "many stakeholders and many variables that are not easily
%     controllable" : vague and uncited. Name two or three concrete actors
%     or constraints instead. Specifics are more persuasive than "many".
% [f] "predictive look at what future land permeability ... looks like" :
%     the forecast itself is authorised (D26/D27), so this claim is fine in
%     substance. But "looks like" is loose for an abstract-adjacent
%     sentence, and the scenario structure (conservation/low-growth,
%     continuation, high-development) is the interesting part. Name it.
% [g] KEEP AND BUILD ON: framing unprotected land as the reason this is a
%     conservation paper rather than a park-management paper is the
%     strongest idea in either file, and it is yours. Make it the spine of
%     Move 1 instead of a middle sentence.

\subsection*{Move 2: What is already modelled, and what is not}
\TODO{Habitat suitability for this species and region is now well covered, including
current and projected suitability with protected-area mismatch. State plainly that
this project does not add another suitability model. Then the pivot: suitability is
not permeability, and no published work reports how permeability between established
cheetah range cores has \emph{changed}. Note that the recent regional multispecies
connectivity assessment excluded cheetah for insufficient data, which is
simultaneously the justification and the constraint.}

\subsection*{Move 3: The methodological frame}
\TODO{Least-cost and circuit-theoretic connectivity \parencite{adriaensen2003leastcost,
mcrae2008circuit}; resistance parameterisation is the dominant source of uncertainty
in this class of model \parencite{zeller2012resistance,zeller2018equal}; hence a
predeclared scenario ensemble rather than a single preferred surface. Graph-level
importance follows \textcite{saura2007pc}; metric choice follows
\textcite{keeley2021metrics}. Design choices are framed against
\textcite{beier2008forks}.}

\subsection*{Move 4 --- Objectives and hypotheses}
\TODO{State the three objectives (change, protection gap, monitoring priority) and
list H1--H4 with their directional predictions. Every hypothesis must be falsifiable
as written; if you cannot describe the result that would refute it in one clause,
rewrite it. See the Hypotheses sheet.}

\paragraph{Terminology.} Throughout, \structconn\ refers to modelled landscape
permeability between fixed range cores. It is not evidence of cheetah movement,
dispersal, or gene flow. Where a location is described as a pinch point, this means
a cell of high modelled current concentration, not an observed crossing.
% This paragraph is CAU-001's mitigation. Keep it in the Introduction, not just
% the Discussion. It cannot be cut for length.
% REVIEW: agreed, and it is the clearest statement of your claim boundary
% anywhere in the manuscript. It must survive every length cut you make.

% =====================================================================
% REVIEW 2026-09-01 . CRAFT PATTERN (the thing to carry into Move 2-4)
%
% Your costliest habit is choosing verbs that claim more than the method
% delivers: "observed" change, cheetahs that "will be roaming", a
% "predictive look" at permeability. Each is a small word doing large
% unearned work. In a paper whose entire contribution rests on separating
% modelled permeability from real movement, these are not style slips,
% they are the failure mode itself.
%
% THE RULE, apply it as you write rather than in revision:
% before every verb describing what you did or what happens, ask whether a
% sensor or a person recorded it, or whether your model produced it. If the
% model produced it, the verb is "modelled", "estimated" or "inferred", and
% the noun takes "modelled" as a modifier. Reserve "observed" for the Dryad
% records and the Limpopo tracking data, which are the only observations in
% this project.
%
% Apply it to your TODO notes too. The scaffold currently contains the same
% error ("observed change in network permeability"), so it is teaching you
% the habit every time you read it.
% =====================================================================
